%% Internal brainstorming manuscript for a possible COR submission.
%% Status of statements:
%%   (i) Croella et al. and the accepted conference paper are prior work;
%%  (ii) AGLR describes the current implementation;
%% (iii) the main experiments are planned;
%%  (iv) the earlier RC2+B+R table is diagnostic evidence and must not be used as a
%%       submission result.
\documentclass[preprint,12pt,authoryear]{elsarticle}

\usepackage[T1]{fontenc}
\usepackage{amsmath}
\usepackage{amssymb}
\usepackage{amsthm}
\usepackage{booktabs}
\usepackage{multirow}
\usepackage{tabularx}
\usepackage{array}
\usepackage{xurl}
\usepackage[linesnumbered,ruled,vlined]{algorithm2e}
\emergencystretch=1em

\newtheorem{proposition}{Proposition}
\newcommand{\method}{AGLR}

\journal{Computers \& Operations Research}
\biboptions{authoryear,round}

\begin{document}

\begin{frontmatter}

%% Working title. Revise after the selection model and main experiments are fixed.
\title{Adaptive Growth-Limited Refinement in MaxSAT-DDD for Exact Train Rescheduling}

\author[vnu]{Tuyen Van Kieu}
\ead{tuyenkv@vnu.edu.vn}

\author[vnu]{Vinh Quang Do}
\ead{25020428@vnu.edu.vn}

\author[vnu]{Khanh Van To\corref{cor1}}
\ead{khanhtv@vnu.edu.vn}
\cortext[cor1]{Corresponding author.}

\affiliation[vnu]{organization={Faculty of Information Technology, VNU University of Engineering and Technology},
                  city={Hanoi},
                  country={Vietnam}}

\begin{abstract}
%% Write this section only after the correctness gate and main experiments have
%% been completed. Do not reuse the accepted conference abstract.
\end{abstract}

\begin{keyword}
train rescheduling \sep dynamic discretization discovery \sep MaxSAT
\sep selective refinement \sep exact optimization
\end{keyword}

\end{frontmatter}

\section{Introduction}
\label{sec:introduction}

Real-time railway control requires a dispatcher to repair a timetable after
delays or temporary capacity losses while preserving route order and resource
safety. Even when the route of every train is fixed, choosing event times
remains difficult because two trains that need incompatible resources must be
ordered. Exact methods must therefore make two kinds of progress within a
short computation time: they must find a feasible schedule and they must
produce a lower bound that certifies its quality
\citep{Cacchiani2014,Lamorgese2018}.

Time-indexed formulations expose the temporal choices explicitly, but a model
containing every admissible time point can be prohibitively large. Dynamic
Discretization Discovery (DDD) addresses this problem by solving a relaxation
on a partial time discretization, checking the resulting schedule in the
original model, and adding time points only when they are needed
\citep{Boland2017}. Croella et al.\ adapted this framework to fixed-route train
rescheduling and used weighted partial MaxSAT as the optimization engine
\citep{Croella2024}.

Our accepted conference paper extended that MaxSAT-DDD implementation with
route-precedence preprocessing and a hybrid encoding for at-most-one
constraints \citep{KieuEtAl2026Conference}. Those components and their reported
experiments constitute prior work. In this manuscript they are used only as a
baseline. The journal study starts from a different bottleneck: after a
restricted MaxSAT solution has been decoded, a DDD iteration may reveal many
travel-time and resource violations. Repairing all of them can enlarge the
Boolean model quickly, whereas repairing too few can require many expensive
MaxSAT calls.

Recent DDD research confirms that refinement size is a substantive algorithmic
choice. Arc-dependent discretizations reduce unnecessary copies in service
network design \citep{VanDykKoenemann2024}; application-specific strategies
have improved refinement and feasible-solution construction in vehicle
scheduling \citep{VanLieshoutVanderSchaft2025}. More recent work selects a small
set of time points that jointly removes a collection of infeasible patterns
\citep{ShuXuBaldacci2026}, or selects fixes that repair all observed conflicts
while minimizing the resulting relaxation \citep{HelberEtAl2026}. These
developments make adding fewer time points insufficient as a novelty claim.
They also expose a more specific question that remains relevant to the MaxSAT
setting: when an iteration can afford only a limited increase in model size,
which conflicts should be repaired immediately and which can safely be
postponed?

This working manuscript investigates the following questions.

\begin{description}
\item[RQ1.] Does limiting model growth during refinement preserve the optimum
returned by the published CSOnet solver?
\item[RQ2.] Does adaptive selection improve the number of instances proved
optimal and PAR-2 relative to CSOnet and the required MaxSAT and MILP baselines?
\item[RQ3.] Is any difference explained by fewer time points and clauses,
fewer DDD iterations, or lower time in conflict handling and MaxSAT?
\item[RQ4.] How sensitive is the method to the delay objective and its three
parameters?
\end{description}

The current implementation provides Adaptive Growth-Limited Refinement (AGLR).
It scans the complete decoded schedule, assigns each conflict a benefit based
on waiting time, recurrence, and shared time points, and selects conflicts by
benefit per unit of immediate model growth. A waiting rule prevents indefinite
postponement, while a rescue step adds all pending conflicts when neither bound
has improved for several iterations. The implementation keeps the published
CSOnet path unchanged and exposes an all-conflicts control through the same new
scan. Section~\ref{sec:method} gives the complete rule used by the code.

\section{Related Work and Positioning}
\label{sec:related}

\subsection{Exact train rescheduling}

Exact train rescheduling has been modeled with continuous-time mixed-integer
programming, time-indexed formulations, constraint programming, and
decomposition. Continuous-time formulations avoid a fixed time grid but often
rely on disjunctive big-$M$ constraints. Time-indexed formulations make resource
choices combinatorial and can give strong relaxations at the price of model
size. Logic-based Benders decomposition exploits the spatial structure of
railway networks \citep{LamorgeseMannino2015,LeutwilerCorman2022}, while
dedicated algorithms target the short response times required in operations
\citep{Sama2017,Bettinelli2017}. Realistic station studies also show that the
choice of formulation and operational scope strongly affects practical
performance \citep{ManninoNakkerud2023}.

The direct baseline for this study is the MaxSAT-DDD model of
\citet{Croella2024}. A restricted interval-assignment problem is encoded as
weighted partial MaxSAT: hard clauses describe temporal feasibility in the
current discretization, and soft clauses encode the delay objective. The
accepted conference study added preprocessing and Boolean-encoding components
to this baseline \citep{KieuEtAl2026Conference}. The present paper does not
claim those components again.

\subsection{Dynamic discretization and refinement}

The original DDD framework alternates between a lower-bounding partial
time-expanded model and a refinement step \citep{Boland2017}. Later work has
made the discretization more expressive, for example by maintaining
arc-dependent rather than node-dependent time sets
\citep{VanDykKoenemann2024}. In a multi-depot vehicle-scheduling application,
\citet{VanLieshoutVanderSchaft2025} separate the construction of the partial
network, duty refinement, feasible-solution recovery, adaptive tolerance, and
post-processing, and evaluate these components individually.

Two recent studies are particularly close to the present question.
\citet{ShuXuBaldacci2026} collect minimal infeasible patterns and solve a joint
refinement problem that removes them with a small set of time points.
\citet{HelberEtAl2026} formalize an optimal refinement problem in which a subset
of candidate fixes must repair all observed conflicts while minimizing
relaxation size. The latter study explicitly separates this task from deciding
which subset of conflicts should be repaired. Our proposed selection model
targets that decision in a MaxSAT-DDD algorithm for train rescheduling.

\subsection{Selection in iterative exact methods}

An iterative exact method can become slower when it retains too much
information. In railway rescheduling, \citet{LeutwilerBonetCorman2023} show that
reusing Benders cuts can reduce the number of iterations but that excessive cut
reuse enlarges the optimization problems and creates a competing slowdown.
Their lazy-selection mechanism motivates measuring both sides of the same
trade-off here: a refinement can prevent a conflict from recurring, but its
time points, variables, and clauses remain in subsequent MaxSAT instances.

\subsection{Positioning}

Table~\ref{tab:positioning} distinguishes the accepted work from the journal
method. AGLR is implemented; its performance remains an empirical question
until the staged campaign has been completed.

\begin{table}[htbp]
\centering
\caption{Positioning of the working journal study. All observed means that the
refinement optimization is required to repair every conflict or pattern passed
to it.}
\label{tab:positioning}
\footnotesize
\begin{tabularx}{\textwidth}{@{}p{30mm}p{25mm}p{45mm}X@{}}
\toprule
Method & Application & Refinement decision & Role in this study \\
\midrule
Croella et al. & Train rescheduling & Add repairs from the current solution &
Original MaxSAT-DDD baseline \\
Accepted conference paper & Train rescheduling & Same refinement layer; new
preprocessing and encoding & Prior-work baseline \\
Shu et al. & Service network design & Jointly minimize points for collected
patterns & Closest joint-refinement work \\
Helber et al. & Service network design & Minimize fixes for all observed
conflicts & Closest optimal-refinement work \\
CSOnet implementation & Train rescheduling & Add conflicts found by its
published local scan & Reproducible internal baseline \\
AGLR & Train rescheduling & Select recurring and shared conflict repairs under
an adaptive growth budget & Implemented journal method under evaluation \\
\bottomrule
\end{tabularx}
\end{table}

\section{Problem Definition and Journal Baseline}
\label{sec:problem}

\subsection{Fixed-route train rescheduling}

Let $\mathcal I$ be the set of trains. Train $i$ follows an ordered sequence of
visits $V_i=(v^i_1,\ldots,v^i_{m_i})$. A visit represents the use of one railway
resource, and $t_v\in\mathbb Z$ is its entry time. For a nonterminal visit $v$,
$s(v)$ denotes its successor and $p_v$ is the minimum time to reach that
successor. The exit time is
\begin{equation}
e_v(t)=
\begin{cases}
t_{s(v)}, & \text{if }s(v)\text{ exists},\\
t_v+p_v, & \text{if }v\text{ is terminal}.
\end{cases}
\label{eq:exit-time}
\end{equation}

Each visit has an earliest entry $\underline t_v$. A subset
$\mathcal A\subseteq \bigcup_i V_i$ also has a target time $a_v$ used by the
delay objective. Let $\mathcal C$ contain pairs of visits whose resources are
incompatible. A schedule is feasible when
\begin{align}
t_v &\ge \underline t_v,
&& v\in \bigcup_i V_i, \label{eq:earliest}\\
t_{s(v)} &\ge t_v+p_v,
&& v\text{ with a successor}, \label{eq:route-precedence}\\
e_v(t)\le t_w\quad &\text{or}\quad e_w(t)\le t_v,
&& \{v,w\}\in\mathcal C. \label{eq:resource-disjunction}
\end{align}
Equation~\eqref{eq:resource-disjunction} treats an occupation interval as
$[t_v,e_v)$, so two trains may touch at an endpoint without overlapping.

For delay $\delta_v(t)=\max\{0,t_v-a_v\}$, the implementation minimizes
\begin{equation}
F(t)=\sum_{v\in\mathcal A} c_v\bigl(\delta_v(t)\bigr).
\label{eq:objective}
\end{equation}
The experiments use three objective families. The finite-step objective
(\texttt{finsteps123}) has a bounded sequence of delay penalties. The infinite
staircase objective (\texttt{infsteps180}) increases by one unit for every
$\Delta=180$ seconds of delay. The continuous objective (\texttt{cont}) charges
one unit per second. This definition follows the current implementation, which
sums costs over visits with a target time; it must be kept synchronized with the
data parser, all solver backends, and the independent validator.

\subsection{Restricted MaxSAT-DDD model}

For each visit $v$, DDD maintains a finite ordered set of time points
$\Lambda_v$. A Boolean literal $d_v(\tau)$ represents $t_v\ge\tau$. Monotonicity
is enforced by
\begin{equation}
d_v(\tau_2)\Rightarrow d_v(\tau_1)
\qquad(\tau_1<\tau_2).
\label{eq:ladder}
\end{equation}
The left endpoint selected in each partial partition gives a relaxed schedule
and a lower bound. Hard clauses encode the temporal information discovered so
far, while weighted soft clauses encode the marginal delay cost
\citep{Croella2024,LiManya2021}. The implementation uses RC2 as its MaxSAT
optimizer \citep{Ignatiev2019}.

Suppose a restricted solution has $t_a=\alpha$ but violates a known precedence
$t_b\ge t_a+h$. Adding the point $\alpha+h$ to $\Lambda_b$ permits the hard
implication
\begin{equation}
\neg d_a(\alpha)\lor d_b(\alpha+h).
\label{eq:refinement-implication}
\end{equation}
Resource violations generate an analogous disjunctive refinement. The loop
terminates when the decoded schedule is feasible in
Equations~\eqref{eq:earliest}--\eqref{eq:resource-disjunction} and its verified
cost equals the lower bound of the restricted MaxSAT model.

\subsection{Conference baseline and new bottleneck}

The accepted conference paper propagated earliest times along each fixed route
before building the initial discretization and used a hybrid pairwise/sequential
counter encoding for suitable at-most-one constraints
\citep{KieuEtAl2026Conference,Sinz2005}. These are fixed characteristics of the
conference baseline in the journal comparison.

The new study changes the refinement layer. Let $\mathcal B_k$ be the set of
violations found after solving iteration $k$. Adding all associated points can
reduce future recurrence but can also increase the number of Boolean variables
and clauses passed to every later RC2 call. Conversely, a small refinement
budget can leave important conflicts unresolved and cause many iterations.
The journal experiments therefore record not only total runtime but also the
number of DDD iterations, detected and selected packages, new time points,
forced orders, variables, clauses, and time spent in encoding, MaxSAT,
validation, and selection.

\section{Adaptive Growth-Limited Refinement}
\label{sec:method}

AGLR changes only the refinement stage of the accepted CSOnet solver. The
restricted MaxSAT model, its objective encoding, and the conditions for an
optimal solution remain unchanged. After each MaxSAT call, AGLR scans every
route arc and every potentially conflicting pair of resource visits. It does
not generate the large collection of bound-implied resource orders used in the
earlier RC2+B prototype.

\subsection{Stable conflict candidates}

A detected candidate $b$ contains a stable conflict key, a type (travel or
resource), and the set $P_b$ of time points required by its refinement clause.
A travel key identifies its two consecutive visits. A resource key stores its
two visits in a fixed order. Current event times are deliberately excluded from
the key, allowing the algorithm to recognize a conflict that reappears at a
later iteration.

For every key, AGLR maintains its number of observations $r_b$ and waiting age
$a_b$. The observation count increases whenever the key is detected. The age
increases when the same key remains pending for another iteration and is reset
when the candidate is selected or disappears. Duplicate time points are
removed before costs are evaluated.

\subsection{Immediate growth and adaptive budget}

Let $D_k$ be the time points present before refinement iteration $k$, and let
$S$ be the candidates already selected in that iteration. The marginal growth
estimate is
\begin{equation}
g(b\mid S)=1+\left|P_b\setminus
\left(D_k\cup\bigcup_{q\in S}P_q\right)\right|.
\label{eq:aglr-growth}
\end{equation}
The unit term represents the primary refinement clause. The set term counts
distinct new time points, including a shared point only once. This inexpensive
quantity guides selection; it is not claimed to equal the final CNF increase,
because ladder and at-most-one encodings may add auxiliary variables and
clauses. The experiments record the actual totals separately.

The estimated growth of applying every candidate is
\begin{equation}
C_{\mathrm{all}}=|\mathcal B_k|+
\left|\bigcup_{b\in\mathcal B_k}P_b\setminus D_k\right|,
\label{eq:aglr-all-growth}
\end{equation}
and the iteration budget is
\begin{equation}
K_k=\max\left\{K_{\min},\left\lceil\rho C_{\mathrm{all}}\right\rceil\right\}.
\label{eq:aglr-budget}
\end{equation}
Thus $K_{\min}$ prevents a very small absolute budget, whereas $\rho$ scales the
budget with the conflicts actually observed in that iteration.

\subsection{Benefit and deterministic selection}
\label{sec:selection-model}

Let $h_b$ be the number of distinct points in $P_b$ that also occur in at least
one other current candidate. AGLR assigns benefit
\begin{equation}
u_b=1+\frac{\min\{a_b,W\}}{W}
+\log_2(1+r_b)+\log_2(1+h_b).
\label{eq:aglr-utility}
\end{equation}
The four terms give every candidate positive value, increase the value of an
older or recurring conflict, and reward points shared by several repairs.
Logarithms limit the influence of a large recurrence or sharing count.

Candidates with $a_b\ge W$ are mandatory and may exceed $K_k$. The remaining
candidates are added one at a time by decreasing
$u_b/g(b\mid S)$, provided that their marginal growth fits the remaining
budget. A fixed ordering of conflict keys resolves ties. If the budget admits
no candidate, the highest-ranked candidate is still selected. If neither the
lower bound nor the best verified feasible value improves for $2W$ iterations,
AGLR selects every pending candidate in a rescue iteration.

\subsection{Complete algorithm and controls}

Algorithm~\ref{alg:aglr} states the implemented loop. The \texttt{legacy}
control retains the original CSOnet conflict handling. The
\texttt{all-scanned} control uses the new complete scan but applies all detected
candidates. The \texttt{adaptive} mode uses Equations
\eqref{eq:aglr-growth}--\eqref{eq:aglr-utility}.

\begin{algorithm}[htbp]
\caption{MaxSAT-DDD with adaptive growth-limited refinement}
\label{alg:aglr}
\footnotesize
\KwIn{Instance, $K_{\min}$, fraction $\rho$, maximum wait $W$}
\KwOut{Verified optimum, or timeout with a verified incumbent and bounds}
Initialize the published CSOnet model and an empty conflict-history table\;
\While{time remains}{
  Solve the current restricted MaxSAT model and update its lower bound\;
  Decode the schedule and receive any feasible heuristic candidate\;
  Discard a heuristic candidate if its objective cannot improve the incumbent
  or if the same schedule has already been checked\;
  Fully verify every remaining heuristic candidate before updating the upper
  bound\;
  Scan all route arcs and resource pairs to create $\mathcal B_k$\;
  \If{$\mathcal B_k$ is empty and the decoded cost equals the lower bound}{
    \Return the decoded schedule as optimal\;
  }
  Update $a_b$, $r_b$, $h_b$, $C_{\mathrm{all}}$, and $K_k$\;
  Select mandatory candidates, then candidates by $u_b/g(b\mid S)$\;
  \If{neither bound improved for $2W$ iterations}{select every pending candidate\;}
  Add the selected time points and refinement clauses, suppressing duplicates\;
}
\Return timeout with the best verified incumbent and final bounds\;
\end{algorithm}

The implementation reports time spent in detection, selection, model
application, MaxSAT, and heuristic validation. It also reports the numbers of
detected, selected, mandatory, and rescued candidates; estimated all-candidate
and selected growth; the maximum waiting age; and the final variable and clause
counts.

\section{Correctness and Progress}
\label{sec:correctness}

AGLR changes when a valid refinement is added, not the meaning of that
refinement. A selected travel candidate adds an implication excluding the
current route-time violation. A selected resource candidate adds the standard
disjunctive refinement excluding the current overlap. Therefore, every added
clause is valid for the original train-rescheduling model.

\begin{proposition}[Per-iteration progress]
If the complete scan detects at least one violation, AGLR selects at least one
candidate and adds its valid refinement.
\end{proposition}
\begin{proof}
Mandatory candidates are selected first. If none is mandatory, the ratio rule
selects every fitting candidate. If none fits, the explicit fallback selects
the highest-ranked candidate even when this exceeds the budget. Its refinement
excludes the detected violating assignment.
\end{proof}

\begin{proposition}[No indefinite postponement]
A conflict key that remains detected in consecutive iterations is selected no
later than the iteration at which its waiting age reaches $W$.
\end{proposition}
\begin{proof}
The history update increases the pending age once per consecutive observation.
The mandatory phase selects every candidate with $a_b\ge W$ without applying
the budget test and then resets its pending age.
\end{proof}

The rescue rule is an additional practical safeguard: it applies every pending
candidate after $2W$ iterations without bound progress. Exactness still follows
from the published DDD termination condition: the solver returns optimal only
when the decoded schedule passes the original feasibility check and its
recomputed cost equals the MaxSAT lower bound. The independent result auditor
rechecks this condition from the stored schedule.

With indexed route arcs and resource-to-visit lists, the complete conflict scan
costs $O(n+m)$ for $n$ visits and $m$ candidate resource pairs. For $B$ detected
candidates, the direct deterministic implementation of repeated ratio selection
is $O(B^2)$ in the worst case. Selection time is recorded so this overhead is
visible rather than absorbed into MaxSAT time.

\section{Computational Study}
\label{sec:experiments}

\subsection{Correctness and reproducibility gate}
\label{sec:correctness-gate}

The verifier now uses the exact rules in
Equations~\eqref{eq:earliest}--\eqref{eq:resource-disjunction}: an event cannot
start before its earliest time, and half-open resource occupations may touch
but may not overlap. Tests cover entry one second before and exactly at the
earliest time, as well as touching intervals and a one-second overlap. Successful
runs also store the complete schedule. A separate Python audit reconstructs the
instance from its text file and checks route times, resource occupations, and
objective value without calling the Rust verifier.

The campaign is accepted only when the following conditions hold:
\begin{enumerate}
\item selector tests cover deterministic ties, shared points, the adaptive
budget, mandatory overflow, and rescue selection;
\item the legacy alias reproduces the objective and optimal status of CSOnet,
and small AGLR instances agree with exhaustive enumeration;
\item every \texttt{status=ok} schedule passes the independent validator, its
cost is recomputed, and its lower and upper bounds agree;
\item JSON metadata match the intended policy, $K_{\min}$, $\rho$, $W$,
timeout, and binary;
\item the commit, source diff, \texttt{Cargo.lock}, dependency patches, compiler,
CPU, memory limits, seeds, data checksums, stdout, and stderr are retained; and
\item tables and figures are generated from raw CSV files rather than copied
manually into this manuscript.
\end{enumerate}

\subsection{Instances and objectives}

The benchmark introduced by \citet{Croella2024} has 24 base instances from two
Norwegian single-track networks. Each base case is provided as an original
instance, a variant with increased track running times, and a variant that also
increases station dwell times. This gives 72 instances for each objective. The
journal campaign uses the complete set for \texttt{finsteps123},
\texttt{infsteps180}, and \texttt{cont}.

The accepted conference results will be included as a prior-work baseline after
their binary, commit, configuration, and CSV provenance have been recovered.
Runtime is reused only when the machine, solver version, limits, and optimality
rule are comparable. The mandatory main-table baselines are the earlier journal
MaxSAT-DDD, CSOnet, Big-M Gurobi, and Big-M CPLEX. A MILP run is counted as
optimal only when it returns a valid schedule and matching lower and upper
bounds within the stated tolerance.

\subsection{Ablation configurations}

Table~\ref{tab:ablation-configurations} defines the three controls implemented
on the CSOnet code base. The original solver name invokes the legacy path. The
AGLR alias invokes either the complete new scan or adaptive selection.

\begin{table}[htbp]
\centering
\caption{Implemented controls for isolating the effect of adaptive selection.}
\label{tab:ablation-configurations}
\begin{tabular}{lll}
\toprule
Configuration & Conflict scan & Selection \\
\midrule
CSOnet legacy & published local scan & published rule \\
AGLR-All & complete new scan & all detected candidates \\
AGLR & complete new scan & Equations~\eqref{eq:aglr-budget}--\eqref{eq:aglr-utility} \\
\bottomrule
\end{tabular}
\end{table}

Parameter selection is separated from final evaluation. The 18 combinations
$K_{\min}\in\{16,32,64\}$, $\rho\in\{0.25,0.5,0.75\}$, and
$W\in\{5,10\}$ first receive one 30-second run on 24 development instances.
The three best configurations receive one 120-second run on the same set. One
configuration is then locked before the complete 72-instance evaluation.

\subsection{Evaluation measures}

Primary measures are the number of valid feasible schedules, the number of
proved optima, timeout and error counts, solved-over-time, and PAR-2. For a
cutoff $T$ and $N$ instances,
\begin{equation}
\mathrm{PAR2}=\frac1N\sum_{i=1}^{N}p_i,\qquad
p_i=
\begin{cases}
t_i,&\text{if instance }i\text{ is completed within }T,\\
2T,&\text{otherwise}.
\end{cases}
\label{eq:par2}
\end{equation}
Runtime medians and speedups are computed on paired instances. A median over a
different solved set is not used to establish that one method is faster.

Mechanism measures include detected and selected candidates, estimated primary
growth, mandatory and rescue selections, maximum age, DDD iterations, variables,
clauses, and time spent in detection, selection, application, MaxSAT, and
heuristic validation. LB and UB traces are retained for timeout cases.

\subsection{Diagnostic pilot}
\label{sec:pilot}

Table~\ref{tab:pilot-comparison} summarizes the already audited comparison that
motivated AGLR. RC2+B+R selected almost every detected constraint when its point
budget was 128. It gave a small gain for \texttt{finsteps123} and regressions for
the other objectives. The results are retained as diagnostic evidence and are
excluded from the main journal comparison.

\begin{table}[htbp]
\centering
\caption{Diagnostic one-run comparison at a 120-second cutoff. PAR-2 includes a
two-cutoff penalty for runs not proved optimal.}
\label{tab:pilot-comparison}
\begin{tabular}{lrrrr}
\toprule
& \multicolumn{2}{c}{RC2-Base} & \multicolumn{2}{c}{RC2+B+R, $K=128$} \\
Objective & Proved & PAR-2 (s) & Proved & PAR-2 (s) \\
\midrule
\texttt{finsteps123} & 72/72 & 1.538 & 72/72 & 1.451 \\
\texttt{infsteps180} & 68/72 & 16.667 & 67/72 & 18.092 \\
\texttt{cont}        & 57/72 & 54.878 & 56/72 & 56.694 \\
\bottomrule
\end{tabular}
\end{table}

At $K=128$, 66 of 72 instances selected all detected constraints, so this run
did not establish a useful selection effect. The bound component also generated
many precedence relations, and heuristic verification dominated some easy
instances. These observations motivated moving the selector directly into
CSOnet, measuring actual model growth, and filtering heuristic candidates before
full validation.

\subsection{Planned result structure}

The final analysis will answer the research questions in the following order:
\begin{enumerate}
\item establish objective equality and schedule validity across configurations;
\item select parameters by proved optima, PAR-2, and estimated selected growth;
\item compare AGLR coverage, PAR-2, and paired runtimes with all mandatory
baselines;
\item use the legacy/all-scanned/adaptive controls to isolate complete scanning
from adaptive selection;
\item explain differences through candidate, clause, iteration, phase-time, and
bound traces; and
\item report objective-specific regressions and failure cases.
\end{enumerate}
No numerical statement will be moved to the abstract or conclusion until the
corresponding raw results and generation script are archived.

\section{Discussion}
\label{sec:discussion}

The current design separates detection from selection. Detection decides which
travel or resource conditions the decoded schedule violates. Selection decides
which valid repairs are added now. The growth budget affects only the latter;
it does not weaken the independent feasibility check or the condition used to
declare an optimum.

Selective refinement can keep the MaxSAT instance compact, while postponement
can add solver iterations. Age and rescue rules control postponement, but the
benefit formula remains a heuristic for computational value. The
all-scanned control and phase telemetry are therefore needed to distinguish a
useful selection effect from overhead introduced by the new scan.

The three objectives may react differently because they create different
numbers and distributions of objective levels. Continuous delay can enlarge the
Boolean objective encoding, whereas the finite-step objective has few levels.
The same refinement policy may therefore move work between conflict handling
and objective optimization in different ways.

External validity remains limited by the fixed-route model and by 72 variants
derived from 24 base cases in two networks. The model does not currently cover
rerouting, cancellation, passenger connections, rolling-stock recovery, or
uncertain running times. A favorable result on the current benchmark should be
reported as evidence for this setting rather than a universal advantage of
budgeted refinement.

\section{Conclusions and Next Steps}
\label{sec:conclusion}

This working manuscript defines AGLR as an implemented extension of the CSOnet
refinement loop. The method combines a complete conflict scan, a deterministic
benefit-to-growth selector, an age rule, and a rescue rule. The code also
filters repeated or non-improving heuristic candidates and records the measures
needed to explain its behavior. The accepted conference contributions remain
prior work and a baseline.

This section is intentionally provisional. The staged parameter search, full
72-instance campaign, mandatory MaxSAT and MILP comparisons, and confirmation
runs must still be completed on the experiment machine. Abstract, numerical
Results, and performance conclusions will be written only from independently
audited outputs that meet the preregistered coverage and PAR-2 thresholds.

\section*{Data and Code Availability}

The evolving implementation, scripts, and benchmark interfaces are maintained
at \url{https://github.com/maxsat-tsp/maxsattrainscheduling}. The submission
artifact will identify the exact commit and provide raw JSON/CSV results,
configuration manifests, data checksums, dependency patches, and scripts that
generate every reported table and figure.

\section*{CRediT Authorship Contribution Statement}

%% Complete only after all authors confirm their roles.

\section*{Declaration of Competing Interest}

%% Replace this comment with the authors' confirmed declaration.

\section*{Acknowledgements}

%% Add only confirmed funding and acknowledgements.

\bibliographystyle{elsarticle-harv}
\bibliography{references}

\end{document}
